\chapter*{Preface}
\addcontentsline{toc}{chapter}{Preface}

In the era of global digital transformation, automatic document understanding and structured information extraction from business documents play a pivotal role across banking, finance, insurance, healthcare, e-commerce, and public administration workflows. Every day, millions of invoices, point-of-sale (POS) receipts, tax forms, contracts, and administrative records are digitized as scanned images or PDF files. However, converting these unstructured or semi-structured document images into structured machine-readable records remains a major bottleneck. This challenge forms the core research scope of \textbf{Document Artificial Intelligence (Document AI)}, and more specifically, the fundamental task of \textbf{Key Information Extraction (KIE)} on \textbf{Visually Rich Documents (VRDs)}.

Unlike plain natural language text, Visually Rich Documents convey semantic meaning not only through lexical words, but also through two-dimensional (2D) spatial layouts, tabular alignments, multi-column regions, and visual typography (font sizes, bolding, borders, and checkboxes). Multi-modal pre-trained Transformers—most notably the LayoutLM family and specifically \textbf{LayoutLMv3}—have achieved remarkable breakthroughs by unifying text tokens, 2D bounding coordinates, and document image patches within a single Transformer backbone. Nevertheless, when fine-tuning standard LayoutLMv3 for sequence labeling on complex KIE benchmarks, three critical architectural limitations persist:
\begin{enumerate}
    \item \textbf{Isotropic 2D coordinate embeddings:} The model projects horizontal coordinates $(x_0, x_1, w)$ and vertical coordinates $(y_0, y_1, h)$ independently using a shared isotropic embedding dimension, failing to decouple intra-line reading continuity from inter-line hierarchical structure.
    \item \textbf{Absence of explicit line/segment structural consensus:} Self-attention operates solely at the subword token level without an explicit intermediate representation of physical text lines (segments), leading to inconsistent IOB2 label predictions within the same semantic text block.
    \item \textbf{Lack of intra-line boundary regularization:} Standard token-level cross-entropy loss treats every token independently and does not explicitly penalize abrupt label transitions between adjacent tokens belonging to the same text line.
\end{enumerate}

Motivated by these challenges, this undergraduate thesis proposes a \textbf{Structure-Aware Multi-Modal LayoutLMv3 Architecture} that integrates three synergistic architectural and optimization enhancements: (1) \textbf{Hierarchical 2D Positional Encoding (HPE)}, (2) a \textbf{Token-Segment Interaction Module (Segment Consensus Module - SCM)} with ReZero-style learnable gating and segment-start embeddings, and (3) an auxiliary \textbf{Intra-Line Boundary Regularization Loss}.

The remainder of this thesis is organized into five main chapters matching the standard academic format:
\begin{itemize}
    \item \textbf{Chapter 1 -- Introduction:} Presents the practical motivation of Document AI and KIE on VRDs, formalizes the multi-modal sequence labeling problem, analyzes existing challenges, and summarizes the main contributions of the thesis.
    \item \textbf{Chapter 2 -- Background and Related Work:} Reviews foundational architectures including Vision-Language Transformers, BIO/IOB2 sequence tagging, the evolution from LayoutLM and LayoutLMv2 to LayoutLMv3, and segment-aware document understanding methods such as StructuralLM, DocFormer, LiLT, GeoLayoutLM, and reading-order aware pipelines.
    \item \textbf{Chapter 3 -- Proposed Methodology:} Formulates the mathematical foundations and algorithmic workflow of the three proposed enhancements (HPE, Token-Segment Interaction with ReZero Gating, and Intra-Line Boundary Regularization) atop the LayoutLMv3 backbone.
    \item \textbf{Chapter 4 -- Experiments and Results:} Details the experimental setup on an NVIDIA A100-SXM4-80GB GPU, the FUNSD and CORD benchmark datasets, entity-level evaluation metrics across 3 random seeds ($\{42, 123, 1993\}$), quantitative comparisons, and systematic ablation studies.
    \item \textbf{Chapter 5 -- Conclusion:} Reserved for final conclusions and future research directions once all base and large ablation suites complete execution.
\end{itemize}

